\chapter{已有能力与缺失环节}

本章要判定的量是：哪些方程和哪些装置行为已经存在，从而后文的步骤必须在它们之外仍然可区分。判定失败的表现是把体积平均、腔温闭环或超温断电再叙述一遍，并称为新的环节。

\section{食品热过程已经具备的物理描述}

\paragraph{要判定的量}
下列陈述是后文的输入，还是后文要重新建立的方程。若是输入，后文只允许在它们的连接处增加判据。

\paragraph{继承的文献层与不继承的层}
每一条都对应第 \ref{ch:lit} 章的一层，这里只列使用方式。

\begin{enumerate}
\item 多孔介质的体积平均给出液、汽、能量的共同骨架 \cite{whitaker1977,datta2007}。蒸比干热更快加热内部，是因为表面冷凝释放潜热；烤形成干壳，是因为蒸发吸热并降低近表面渗透率。本文不把这个差异写成两种互不相关的设备模型。
\item 蛋白变性、淀粉糊化、美拉德褐变的速率可用指数温度因子乘以转化率的幂，再乘以含水开关 \cite{vanboekel2008}。品质是反应进度，不是到达某一温度的时刻。
\item 热力安全进度沿冷点温度历史积分 \cite{holdsworth1997}。冷点随边界移动，不一定是几何中心 \cite{feyissa2013}。
\item 由表面温度反演内部热流是病态的，高频噪声随厚度被放大 \cite{ozisik2000}。
\item 物理信息网络和神经算子可以加速参数化演化 \cite{raissi2019,li2021fno,lu2021deeponet}。它们的外推范围由训练轨迹里的执行器覆盖决定。本次没有训练这些算子。
\item 结构因果模型里的干预，是替换结构方程的输入后再积分；条件相关保留控制策略本身 \cite{pearl2009}。
\item 模型预测控制给出有限时域的执行量 \cite{rawlings2017}。控制屏障把该执行量投影到由测量定义的安全集 \cite{ames2017cbf}。本次日志里的加热律不是这个优化问题的解。
\end{enumerate}

\paragraph{缺失的环节}
缺失的不是上列方程，而是这些模块之间的判据：何时后验已经够用，何时必须改变边界来买信息，信息写入哪几个参数，效应沿哪条路径改变品质，以及品质模型失效时谁接管执行器。后文的命题都是这些判据的可否定形式。

\paragraph{本节不能推出的结论}
“文献里已有方程”不能推出那些方程的数值解已经在本文的日志里。规格与十六条控制体的差别写在附录 \ref{app:discrete}。

\section{现有装置实际闭合了哪一段}

\paragraph{要判定的量}
一条现有路径闭合的是哪一个量，哪一些决策量没有进入后验。闭合腔温并不闭合冷点。

\paragraph{继承的文献层与不继承的层}
下表不引用某一商品的控制逻辑。它只区分已经能够稳定调节的边界量和仍然隐藏的食品量。不继承“蒸烤一体等于状态估计已经合并”这一说法：合并的是执行器，不是内部状态。

\begin{table}[ht]
\centering
\caption{已有路径的信息止点}
\begin{tabularx}{\textwidth}{lXX}
\toprule
路径 & 已闭合的量 & 未进入后验的量 \\
\midrule
机械定时 & 无 & 全部内部状态 \\
腔温闭环 & 空气温度一点 & 冷点、含水、进度 \\
按重量调用曲线 & 申报质量与类别 & 初温、厚度、脂肪、形状 \\
穿刺探针 & 刺入点温度 & 未刺中的冷点；刺入改变局部场 \\
表面红外 & 辐射亮度 & 发射率、反射、内部梯度 \\
腔体湿度闭环 & 相对湿度设定误差 & 食品内部水分与冷凝位置 \\
图像分类后调用曲线 & 外观类别 & 中心进度；蒸汽遮挡使图像中断 \\
\bottomrule
\end{tabularx}
\end{table}

机械定时的信息止点最干净：它不读取任何食品量，停止时刻是开环函数。本次核的开环策略就是这一类，菜单时间 \(t=900+18000(L-0.015)\) 秒。15 mm 为 900 s，30 mm 为 1170 s，50 mm 为 1530 s。交叉表里开环的中位 \(e_{\alpha}=-0.843\)，提前比例 \(0.667\)：薄件被推到进度上沿，厚件停在设计线之下。这是“曲线按厚度档标定、内部热阻从未被这次过程观察”的数值形式，不是对某一商品菜单的测定。

腔温闭环在核里是 \(u=\mathrm{clip}(0.04(160-y_c)+0.6,0,1)\)。它闭合空气温度，不估计冷点，也不停机。中位能量 3310、中位 \(e_{\alpha}=+0.150\)，是时域终点上进度饱和，不是推荐的成功。

识别食材后调用固定曲线，曲线不随这一次的后验更新。图像在蒸汽遮挡时中断，所以外观类别不能充当冷点进度的连续观测。

\paragraph{本节不能推出的结论}
表中“未进入后验”不等于那些传感器不存在。它等于那些读数没有被写成 \(\phi\) 的充分统计量。湿度闭环最小化湿度误差，与后文最大化参数信息，可以驱动同一阀门。

\section{不构成新增环节的做法}

\paragraph{要判定的量}
一条步骤是否与下列六条可区分。不可区分就不写入新增环节。

\paragraph{六条}
\begin{enumerate}
\item 按重量或图像选择时间—温度曲线。
\item 穿刺温度到达阈值后停止加热。
\item 用腔温误差调节加热器，用湿度误差调节蒸汽。
\item 用回归模型把腔温历史映射成剩余时间。
\item 漏电、超温、门锁各自独立动作，与品质估计无关。
\item 把传感器曲线映射到三维着色作为展示。
\end{enumerate}

第 4 条可以在观测分布上误差很小，同时在干预下失败：蒸汽加大可能只是控制器对当前低温的反应。相关保留了这条策略边。第 5 条是必要的硬保护，但它不读取品质，也不证明品质估计是对的。第 6 条不更新内部热阻，也不改变停止判据。演示章里的着色若不能指回状态分量，就属于这一条。

\paragraph{本节不能推出的结论}
把这六条排除，不能自动推出后文步骤已经实现。排除的是“声称”，实现以日志和断言为准。

\section{新增环节}

\paragraph{要判定的量}
每一条新增环节是否有输入、输出和落点。没有落点的句子不进入本章。

\begin{table}[ht]
\centering
\caption{加入闭环的环节}
\begin{tabularx}{\textwidth}{lXX}
\toprule
环节 & 输入与输出 & 落点 \\
\midrule
决策状态 & 全场压缩为冷点、表面、进度、热阻 & \eqref{eq:z} \\
被动估计上限 & 稳态观测不含内部热阻 & 命题 \ref{prop:steady} \\
信息动作 & 相变变号之后接对流 & 命题 \ref{prop:pulse} \\
路径效应 & 同一执行器的冷点项与表面项分开 & \eqref{eq:path} \\
有界更新 & 残差越界则丢弃修正项 & \eqref{eq:bnp} \\
测量隔离投影 & 安全不等式不含品质估计 & \eqref{eq:cbf} \\
\bottomrule
\end{tabularx}
\end{table}

这六条都是连接，不是新的本构、新的熵或新的屏障定义。连接若与已发表的技术方案相同，改为引用并退出创新表述。本次阅读没有遇到完全相同的四条连接，未读到不等于不存在。

已运行核与这六条的距离需要逐条写明，避免把规格落点写成已执行。

\begin{itemize}
\item 决策状态在日志摘要里只保留停机进度、能量、热阻方差和终点褐变。温度离散、表面液相和结壳没有单独的结果列。
\item 被动上限有断言：无相变平衡时温差极差为 0。
\item 信息动作在核里是 \([150,190)\) 秒的固定窗，不是探针库上的打分。断言只比较了两段脉冲与单独蒸汽的滞后，差值 \(40.846-38.912=1.934>1\)。
\item 路径效应没有在评分核里按链式法则积分。
\item 有界残差网络没有训练，核里没有 \(\vect{r}_{\psi}\)。
\item 测量隔离在断言里是 \(210^\circ\mathrm{C}\) 的测试函数；在物理步进里是 \(230^\circ\mathrm{C}\) 以上的削减，而加热器平衡在 \(196.5^\circ\mathrm{C}\)，削减不动作。
\end{itemize}

\paragraph{本节不能推出的结论}
表中的落点是后文必须证明或实现的位置。落点存在，不等于该位置的数值结果已经产生。
